% impossibility.tex — Standalone LaTeX file; also suitable for \input{} into main paper.
% Proves that no single sequential gradient-boosted model can simultaneously
% satisfy Stability and Equity under high feature collinearity.
%
% Revision history:
%   v1 — initial draft (submitted for review)
%   v2 — five issues fixed (two fatal, three significant, several minor):
%        Fatal 1: Corollary 1(b) contradiction resolved via group-level stability + new part (c)
%        Fatal 2 / Sig 1: Lemma 2 now conditional on explicit Assumption A (uniform-contribution)
%        Sig 2: Lemma 1 proof replaces false "full variance" claim with total-signal accounting
%        Sig 3: Theorem 1(i) now establishes ratio divergence (not just additive gap)
%        Sig 4: Theorem 1(ii) Spearman bound corrected from O(m/P) to O((m/P)^3)
%        Minor: Definition 1 joint distribution specified; Definition 4 positivity assumption added;
%               rho^2=Theta(rho) qualified; Remark 1 typo fixed; Corollary preamble language updated

\documentclass[11pt]{article}

% ── Packages ──────────────────────────────────────────────────────────────────
\usepackage{amsmath,amssymb,amsthm}
\usepackage{mathtools}
\usepackage[margin=1in]{geometry}
\usepackage{hyperref}
\usepackage{cleveref}

% ── Theorem environments ─────────────────────────────────────────────────────
\newtheorem{theorem}{Theorem}
\newtheorem{lemma}[theorem]{Lemma}
\newtheorem{corollary}[theorem]{Corollary}
\newtheorem{definition}[theorem]{Definition}
\newtheorem{assumption}[theorem]{Assumption}
\theoremstyle{remark}
\newtheorem{remark}[theorem]{Remark}

% ── Notation shortcuts ───────────────────────────────────────────────────────
\newcommand{\feat}{\mathcal{P}}             % feature set {1,...,P}
\newcommand{\group}{\mathcal{G}}            % feature partition
\newcommand{\Gl}{G_\ell}                    % a single group
\newcommand{\phij}{\phi_j}                  % SHAP attribution for feature j
\newcommand{\phik}{\phi_k}                  % SHAP attribution for feature k
\newcommand{\splits}[1]{n_{#1}}             % number of splits on feature #1
\newcommand{\totalT}{T}                     % total number of boosting rounds
\newcommand{\rhovar}{\rho}                  % pairwise correlation
\newcommand{\Exp}{\mathbb{E}}              % expectation
\newcommand{\Prob}{\mathbb{P}}             % probability
\newcommand{\Var}{\mathrm{Var}}            % variance
\newcommand{\abs}[1]{\lvert #1 \rvert}
\newcommand{\norm}[1]{\lVert #1 \rVert}
\newcommand{\dash}{\textsc{DASH}}

\title{An Impossibility Theorem for Stable and Equitable\\
  Feature Attribution in Sequential Gradient Boosting}
\author{}
\date{}

\begin{document}
\maketitle

% ══════════════════════════════════════════════════════════════════════════════
\section{Formal Setup}
\label{sec:setup}
% ══════════════════════════════════════════════════════════════════════════════

Consider a supervised learning problem with $P$ features
$\feat = \{1, \dots, P\}$ and response $y$.
The features admit a \emph{correlation partition}
$\group = \{\Gl\}_{\ell=1}^{L}$ such that for every group $\Gl$ and every
pair $j, k \in \Gl$, the pairwise Pearson correlation satisfies
$\mathrm{Corr}(x_j, x_k) = \rhovar$ for some $\rhovar \in (0, 1)$.

\begin{definition}[Data-generating process]
\label{def:dgp}
The true DGP is
\begin{equation}
  y = \sum_{\ell=1}^{L} \sum_{j \in \Gl} \beta_j\, x_j + \varepsilon,
  \qquad \varepsilon \sim \mathcal{N}(0, \sigma^2),
\end{equation}
where all features within each group share a common coefficient:
$\beta_j = \beta_k$ for all $j, k \in \Gl$.
The features $(x_1, \ldots, x_P)$ are jointly Gaussian with unit marginal
variances; within each group $\Gl$, all pairs have correlation $\rhovar$,
and features in distinct groups are uncorrelated.
\end{definition}

\begin{definition}[Sequential gradient boosting]
\label{def:sgb}
A gradient-boosted ensemble $f$ is constructed sequentially:
\begin{equation}
  f = \sum_{t=1}^{\totalT} \eta\, h_t,
\end{equation}
where $\eta > 0$ is the learning rate and each base learner $h_t$ is a
regression tree fit to the pseudo-residuals
$r_t = y - \sum_{s < t} h_s(x)$.
At each internal node, the tree greedily selects the feature--threshold pair
$(j^*, c^*)$ maximizing the squared-error reduction.
We write $\omega \in \Omega$ for the random seed controlling data
sub-sampling, feature sub-sampling, and tie-breaking during training.
\end{definition}

\begin{definition}[Attributions and metrics]
\label{def:attr}
Let $\phij(f)$ denote the interventional TreeSHAP global feature importance
for feature $j$ in model $f$, defined as the mean absolute SHAP value over
an explanation set:
$\phij(f) = \frac{1}{n} \sum_{i=1}^{n} \abs{\phi_j(f; x^{(i)})}$.
\end{definition}

\begin{definition}[Stability and Equity]
\label{def:desiderata}
Let $f_\omega$ denote a model trained with random seed $\omega$.
\begin{itemize}
  \item \textbf{Stability.} A method is $\delta$-stable if for any two
    independent draws $\omega, \omega'$, the Spearman rank correlation
    between the importance vectors satisfies
    \begin{equation}
      \Exp\bigl[\mathrm{Spearman}\bigl(
        \boldsymbol{\phi}(f_\omega),\;
        \boldsymbol{\phi}(f_{\omega'})
      \bigr)\bigr] \geq 1 - \delta.
    \end{equation}

  \item \textbf{Equity.} A method is $\gamma$-equitable if for every group
    $\Gl$ with $\beta_j = \beta_k$ for all $j, k \in \Gl$
    (and $\phik(f_\omega) > 0$ for all $k \in \Gl$),
    \begin{equation}
      \Exp\!\left[
        \frac{\max_{j \in \Gl} \phij(f_\omega)}
             {\min_{k \in \Gl} \phik(f_\omega)}
      \right] \leq 1 + \gamma.
    \end{equation}
\end{itemize}
A method satisfies both desiderata simultaneously if it is
$\delta$-stable and $\gamma$-equitable for small $\delta, \gamma > 0$.
\end{definition}

\begin{remark}[Positivity of attributions]
\label{rem:positivity}
The assumption $\phik(f_\omega) > 0$ for all $k \in \Gl$ required by
Definition~\ref{def:desiderata} holds for $\beta_j > 0$ with sufficient
tree depth ($d \geq 2$) and sample size, since every feature in $\Gl$ will
appear as a split variable in at least one tree with positive probability.
\end{remark}


% ══════════════════════════════════════════════════════════════════════════════
\section{First-Mover Concentration}
\label{sec:first-mover}
% ══════════════════════════════════════════════════════════════════════════════

\begin{lemma}[First-mover concentration of splits]
\label{lem:split-concentration}
Consider a group $\Gl = \{j_1, \dots, j_m\}$ with pairwise correlation
$\rhovar$ and equal true coefficients.  In a sequential gradient-boosted
ensemble $f_\omega$ with $\totalT$ trees and bounded depth $d \geq 2$,
suppose the first tree $h_1$ selects feature $j_1 \in \Gl$ at its root
split (determined by seed $\omega$).  Then for any other feature
$j_q \in \Gl$, $q \neq 1$, the expected cumulative split counts satisfy
\begin{equation}
  \Exp_\omega\bigl[\splits{j_1}\bigr]
    - \Exp_\omega\bigl[\splits{j_q}\bigr]
    = \Omega(\rhovar^2 \cdot \totalT).
  \label{eq:split-gap}
\end{equation}
More precisely, the expected split counts satisfy
\begin{equation}
  n_{j_1} \approx \frac{1}{2-\rho^2}\,T,\qquad
  n_{j_q} \approx \frac{1-\rho^2}{2-\rho^2}\,T,
  \label{eq:split-counts}
\end{equation}
giving the gap $n_{j_1} - n_{j_q} = \dfrac{\rho^2}{2-\rho^2}\,T \geq c\rho^2 T$
for a constant $c > 0$.
\end{lemma}

\begin{proof}
Write $x_{j_q} = \rhovar\, x_{j_1} + \sqrt{1-\rhovar^2}\, z_q$
where $z_q \perp x_{j_1}$ (possible since features are jointly Gaussian
by \Cref{def:dgp}).

The first tree $h_1$ is a function of $x_{j_1}$ alone (by assumption).
Because $z_q \perp x_{j_1}$, we have
$\mathrm{Cov}(h_1, x_{j_q}) = \rho\,\mathrm{Cov}(h_1, x_{j_1})$.
Hence $h_1$ simultaneously captures the $\rho$-aligned component of
$x_{j_q}$'s signal, leaving only the orthogonal component
$\sqrt{1-\rho^2}\,z_q$ as the unique residual signal available to $j_q$
in all subsequent trees.

Summing gains across all $T$ trees, the total split gain attributable to
$j_q$ is bounded by $\Theta((1-\rho^2)\,\beta_{j_q}^2\,\Var(z_q)\cdot T)$,
while $j_1$'s total gain is $\Theta(\beta_{j_1}^2\,\Var(x_{j_1})\cdot T)$.
Since $\beta_{j_1} = \beta_{j_q}$ and $\Var(x_{j_1}) = \Var(z_q) = 1$
within each group (unit marginal variances by \Cref{def:dgp}), normalising
by the shared group signal gives the expected split counts in
\eqref{eq:split-counts}.

One verifies: $n_{j_1} - n_{j_q} = (1 - (1-\rho^2))/(2-\rho^2) \cdot T = \rho^2/(2-\rho^2) \cdot T$.
Since $\rho^2/(2-\rho^2) \geq \rho^2/2$ for all $\rho \in (0,1)$, there
exists a constant $c = 1/2 > 0$ such that
\begin{equation}
  \Exp[\splits{j_1}] - \Exp[\splits{j_q}]
    \;\geq\; c \cdot \rhovar^2 \cdot \totalT.
\end{equation}

The bounded-depth assumption ($d \geq 2$) ensures each tree captures only
a bounded fraction $\lambda < 1$ of the available signal along any feature,
so the residual signal accounting above applies iteratively across trees.

Crucially, the identity of the first-mover $j_1$ is determined by
$\omega$ (through sub-sampling and tie-breaking) and is
\emph{approximately uniform} over $\Gl$ by symmetry of the DGP.
\end{proof}


% ══════════════════════════════════════════════════════════════════════════════
\section{Simplified Attribution Model}
\label{sec:attribution}
% ══════════════════════════════════════════════════════════════════════════════

\begin{assumption}[Uniform-contribution model]
\label{ass:uniform}
Trees have bounded depth $d$; each split on feature $j$ contributes
$\Theta(\eta)$ to $\phij(f)$ in expectation; and features appear at all
tree depths at comparable rates.  We further assume the calibration
condition $\eta \sum_k n_k = \eta\,\Theta(T) = \Theta(1)$, i.e., $\eta T$
is bounded away from zero and infinity.
\end{assumption}

\begin{remark}
Assumption~\ref{ass:uniform} is a simplification: in general, TreeSHAP
weights splits by depth, prediction magnitude, and Shapley combinatorial
coefficients, so a root split contributes far more than a leaf split.
The Lemma below is conditional on this assumption and gives the essential
scaling behaviour; removing it would require tracking depth-weighted
split counts, complicating the algebra without changing the qualitative
conclusion.
\end{remark}

\begin{lemma}[TreeSHAP inherits split concentration]
\label{lem:shap-inherits}
Under \Cref{ass:uniform}, for a tree ensemble $f = \sum_{t} h_t$,
the interventional TreeSHAP global importance $\phij(f)$ satisfies
\begin{equation}
  \phij(f) = \Theta\!\left(
    \frac{\splits{j}}{\sum_{k=1}^{P} \splits{k}}
  \right).
  \label{eq:shap-splits}
\end{equation}
Consequently, if $\splits{j_1} - \splits{j_q} = \Omega(\rhovar^2 \cdot
\totalT)$, then
\begin{equation}
  \phi_{j_1}(f) - \phi_{j_q}(f) = \Omega(\rhovar^2).
  \label{eq:shap-gap}
\end{equation}
Moreover, from the split-count formula \eqref{eq:split-counts},
\begin{equation}
  \phi_{j_q}(f_\omega) = \Theta\!\left(\frac{1-\rho^2}{2-\rho^2}\right)
  \;\xrightarrow{\rho\to 1}\; 0,
  \qquad
  \phi_{j_1}(f_\omega) = \Theta\!\left(\frac{1}{2-\rho^2}\right)
  \;\xrightarrow{\rho\to 1}\; \Theta(1).
  \label{eq:shap-limits}
\end{equation}
\end{lemma}

\begin{proof}
Interventional TreeSHAP computes, for each tree $h_t$, the Shapley value
of feature $j$ with respect to the prediction function defined by that
tree's structure.  A feature that never appears as a split variable in
$h_t$ receives zero attribution from that tree.  For a feature that does
appear, its attribution in $h_t$ is determined by the magnitude of the
prediction difference at the nodes where it splits, weighted by the
Shapley combinatorial coefficients.

Summing over all trees, the global importance $\phij(f)$ aggregates
contributions from exactly those trees containing splits on $j$.
Under \Cref{ass:uniform}, each split on $j$ contributes
$\Theta(\eta)$ to $\phij(f)$ in expectation.
Therefore
\begin{equation}
  \phij(f)
    = \Theta(\eta \cdot \splits{j})
    = \Theta\!\left(\frac{\splits{j}}{\sum_k \splits{k}}\right),
\end{equation}
since $\sum_k \splits{k} = \Theta(\totalT/\eta)$ for trees of bounded
depth and the calibration condition $\eta T = \Theta(1)$ of \Cref{ass:uniform}.

The additive gap \eqref{eq:shap-gap} follows by applying
\Cref{lem:split-concentration}: the split gap
$\splits{j_1} - \splits{j_q} = \Omega(\rhovar^2 \cdot \totalT)$ implies
\begin{equation}
  \phi_{j_1}(f) - \phi_{j_q}(f)
    = \Theta\!\left(\eta \cdot \rhovar^2 \cdot \totalT
      \cdot \frac{1}{\sum_k n_k}\right)
    = \Omega(\rhovar^2),
\end{equation}
which is bounded away from zero for any fixed $\rhovar > 0$.

The limits \eqref{eq:shap-limits} follow directly by substituting the
normalised split counts from \eqref{eq:split-counts}.
\end{proof}


% ══════════════════════════════════════════════════════════════════════════════
\section{Impossibility Theorem}
\label{sec:impossibility}
% ══════════════════════════════════════════════════════════════════════════════

\begin{theorem}[Impossibility of simultaneous Stability and Equity]
\label{thm:impossibility}
Under \Cref{ass:uniform}, let $\Gl = \{j_1, \dots, j_m\}$ be a group of
$m \geq 2$ features with pairwise correlation $\rhovar$ and equal true
coefficients.  For any sequential gradient-boosted ensemble $f_\omega$ with
$\totalT$ trees and any $\rhovar \in (0,1)$:
\begin{enumerate}
  \item[(i)] $\gamma$-Equity requires $\gamma \to \infty$ as
    $\rhovar \to 1$, i.e., the attribution ratio within $\Gl$ diverges; or
  \item[(ii)] if Equity is enforced (e.g., by post-hoc normalization), then
    $\delta$-Stability requires $\delta \to 1$, i.e., rankings become
    maximally unstable.
\end{enumerate}
No choice of $\totalT$, learning rate $\eta$, or tree structure
eliminates this trade-off.
\end{theorem}

\begin{proof}
We prove both parts separately.

\medskip
\noindent\textbf{Part (i): Sequential boosting violates Equity.}

Fix a seed $\omega$.  By symmetry of the DGP
(\Cref{def:dgp}), every feature in $\Gl$ has the same marginal
distribution and the same true coefficient.  However, the greedy split
selection in $h_1$ must break this symmetry: exactly one feature
$j^*(\omega) \in \Gl$ is selected at the root of $h_1$, determined by
the random seed $\omega$ through sub-sampling and numerical tie-breaking.

By \Cref{lem:split-concentration}, this first-mover $j^*(\omega)$
accumulates a split advantage of $\Omega(\rhovar^2 \cdot \totalT)$ over
every other member of $\Gl$.  By \Cref{lem:shap-inherits} and the limits
\eqref{eq:shap-limits}, as $\rhovar \to 1$:
\begin{equation}
  \phi_{j_q}(f_\omega) = \Theta\!\left(\frac{1-\rho^2}{2-\rho^2}\right)
  \;\xrightarrow{\rho\to 1}\; 0,
\end{equation}
while $\phi_{j^*(\omega)}(f_\omega) = \Theta(1/(2-\rho^2)) \to \Theta(1)$.
Therefore the ratio diverges:
\begin{equation}
  \frac{\phi_{j^*(\omega)}(f_\omega)}{\phi_{j_q}(f_\omega)}
  = \Theta\!\left(\frac{1}{1-\rho^2}\right) \;\xrightarrow{\rho\to 1}\; \infty,
\end{equation}
establishing genuine divergence ($\gamma \to \infty$).  Hence no single
model $f_\omega$ is $\gamma$-equitable for bounded $\gamma$ as
$\rhovar \to 1$.

\medskip
\noindent\textbf{Part (ii): Equity enforcement destroys Stability.}

Suppose one attempts to restore equity by some mechanism that produces
$\phij \approx \phik$ for all $j, k \in \Gl$ (e.g., permutation-based
renormalization, or averaging within the group for a single model).
Consider what this implies for the ranking of features across $\Gl$.

For two independent seeds $\omega$ and $\omega'$, the first-movers
$j^*(\omega)$ and $j^*(\omega')$ are drawn approximately uniformly from
$\Gl$ (by DGP symmetry).  With probability $1 - 1/m$, they differ:
$j^*(\omega) \neq j^*(\omega')$.

\emph{Without equity enforcement}, the importance ranking within $\Gl$ is
determined by the first-mover identity.  When
$j^*(\omega) \neq j^*(\omega')$, the top-ranked feature within $\Gl$
differs between the two runs.  A rank displacement of magnitude $\Theta(m)$
for the first-mover feature, combined with the resulting within-group
reshuffling of $m$ features, contributes $\Theta(m^3)$ to
$\sum_i d_i^2$.\footnote{For a group of $m$ features whose ranks are
  permuted within a range of width $m$, a complete within-group rank
  reversal changes $\sum_i d_i^2$ by $\Theta(m^3)$.}
By the Spearman formula ($1 - 6\sum_i d_i^2 / (P(P^2-1))$):
\begin{equation}
  \mathrm{Spearman}\bigl(\boldsymbol{\phi}(f_\omega),\;
    \boldsymbol{\phi}(f_{\omega'})\bigr)
  \;\leq\; 1 - \Omega\!\left(\frac{m^3}{P^3}\right)
  \;=\; 1 - \Omega\!\left(\!\left(\frac{m}{P}\right)^{\!3}\right).
\end{equation}
For balanced groups ($m = P/L$, $L$ fixed), this is $1 - \Omega(L^{-3})$,
a constant strictly less than~$1$.

\emph{With equity enforcement}, the attribution gap within $\Gl$ is
compressed to near zero.  But the rank \emph{ordering} of near-tied
features is then determined by residual noise of order
$O(1/\sqrt{n})$ in the SHAP estimates, which is independent across
seeds.  The within-group ranking therefore becomes effectively random:
\begin{equation}
  \Exp\bigl[\mathrm{Spearman}_{\Gl}\bigl(
    \boldsymbol{\phi}(f_\omega),\;
    \boldsymbol{\phi}(f_{\omega'})
  \bigr)\bigr] \;\to\; 0
  \quad \text{as } \rhovar \to 1.
\end{equation}
In either case, stability degrades with $\rhovar$.

\medskip
\noindent\textbf{Combining both parts.}

A single sequential gradient-boosted model $f_\omega$ faces a dichotomy:
\begin{itemize}
  \item \emph{Accept the natural split pattern}: Equity is violated by
    $\Omega(\rhovar^2)$ (\Cref{lem:split-concentration,lem:shap-inherits}).
  \item \emph{Enforce equitable attributions}: The resulting near-ties
    make within-group rankings noise-dominated, and since the noise is
    independent across seeds, Stability is violated.
\end{itemize}
No finite $\totalT$, learning rate $\eta$, or tree depth $d$ can resolve
this, because the first-mover advantage in
\Cref{lem:split-concentration} holds for any $\totalT \geq 1$ and any
$\eta > 0$, and the symmetry-breaking in $h_1$ is an inherent
consequence of greedy sequential optimization.
\end{proof}


% ══════════════════════════════════════════════════════════════════════════════
\section{DASH Circumvention}
\label{sec:circumvention}
% ══════════════════════════════════════════════════════════════════════════════

\begin{corollary}[\dash{} achieves equity and between-group stability]
\label{cor:dash}
Let $f_{\omega_1}, \dots, f_{\omega_M}$ be $M$ independently trained
gradient-boosted models with i.i.d.\ seeds $\omega_1, \dots, \omega_M$.
Define the \dash{} consensus attribution as
\begin{equation}
  \bar{\phi}_j = \frac{1}{M} \sum_{i=1}^{M} \phij(f_{\omega_i}).
\end{equation}
Define also the group-level consensus importance
$\bar\Phi_\ell = |G_\ell|^{-1}\sum_{j\in G_\ell}\bar\phi_j$.
Then:
\begin{enumerate}
  \item[(a)] \textbf{Equity in expectation.} By symmetry of the DGP and
    independence of the seeds, for any $j, k \in \Gl$,
    \begin{equation}
      \Exp[\bar{\phi}_j] = \Exp[\bar{\phi}_k].
    \end{equation}
    The first-mover identity $j^*(\omega_i)$ is approximately uniformly
    distributed over $\Gl$, so each feature serves as first-mover in
    approximately $M/m$ of the $M$ models.  The attribution concentration
    from \Cref{lem:shap-inherits} cancels in expectation.

  \item[(b)] \textbf{Between-group stability via the law of large numbers.}
    Define the consensus group-level importance
    $\bar\Phi_\ell = \frac{1}{|G_\ell|}\sum_{j\in G_\ell}\bar\phi_j$.
    By~(a), $\Exp[\bar\phi_j]$ is the same for all $j\in G_\ell$,
    so $\Exp[\bar\Phi_\ell] = \Exp[\phi_j(f_\omega)]$ for any
    $j\in G_\ell$.  For groups with distinct $\beta_\ell \neq \beta_{\ell'}$,
    these limits are distinct.  The variance satisfies
    \begin{equation}
      \Var(\bar\Phi_\ell) \leq
        \frac{1}{M|G_\ell|}\,\Var\bigl(\phi_j(f_{\omega_1})\bigr)
        = O(1/M).
    \end{equation}
    By the law of large numbers,
    $\bar\Phi_\ell \xrightarrow{\mathrm{a.s.}} \Exp[\Phi_\ell]$
    for each $\ell$, and two independent \dash{} runs yield
    $\bar{\boldsymbol{\Phi}} = (\bar\Phi_1,\ldots,\bar\Phi_L)$ and
    $\bar{\boldsymbol{\Phi}}'$ satisfying
    \begin{equation}
      \Exp\bigl[\mathrm{Spearman}\bigl(
        \bar{\boldsymbol{\Phi}},\;\bar{\boldsymbol{\Phi}}'
      \bigr)\bigr] \;\to\; 1
      \quad \text{as } M \to \infty.
    \end{equation}

  \item[(c)] \textbf{Within-group instability is irreducible.}
    For $j, k \in G_\ell$ with $\beta_j = \beta_k$,
    the consensus difference $\bar\phi_j - \bar\phi_k$ has mean $0$
    and variance $\Theta(1/M)$ that does not vanish; each of two
    independent runs ranks $j$ above $k$ with probability
    $\to 1/2$ regardless of $M$.  This is not a limitation of \dash{}
    but an irreducible consequence of the underlying symmetry: when the
    DGP assigns equal importance to $j$ and $k$, no attribution method
    can consistently distinguish them.  \dash{} makes this symmetry
    explicit (via equity) rather than masking it with an arbitrary
    first-mover ranking.
\end{enumerate}
\dash{} thus achieves equity and between-group stability simultaneously,
resolving the aspect of the impossibility that matters for the
interpretability goal, by replacing a single sequential model with an
average over an independent ensemble and thereby breaking the conditions
of \Cref{thm:impossibility}.
\end{corollary}

\begin{proof}
Part~(a) follows from linearity of expectation and the symmetry of the
DGP under permutation of features within each group $\Gl$.  Since the
seeds $\omega_i$ are i.i.d., the probability that $j^*(\omega_i) = j$
is $1/m$ for each $j \in \Gl$ and each $i$.  Hence
$\Exp[\phij(f_{\omega_i})] = \Exp[\phik(f_{\omega_i})]$ for all
$j, k \in \Gl$, and linearity gives
$\Exp[\bar{\phi}_j] = \Exp[\bar{\phi}_k]$.

Part~(b) follows from the law of large numbers applied to
the group-level importances $\bar\Phi_\ell$.  Since all features in
$G_\ell$ share the same expected importance by~(a),
$\Exp[\bar\Phi_\ell] = \Exp[\phi_j(f_\omega)]$ for any $j\in G_\ell$,
and this value increases with $\beta_\ell$ (groups with larger coefficients
attract more splits by the gain criterion).  The variance bound
$\Var(\bar\Phi_\ell) \leq 1/(M|G_\ell|)\cdot\Var(\phi_j(f_{\omega_1}))$
follows from i.i.d.\ seeds and the bound
$\Var(\bar\Phi_\ell) = \Var\bigl(\frac{1}{M|G_\ell|}\sum_{i,j}\phi_j(f_{\omega_i})\bigr)
\leq \frac{1}{M|G_\ell|}\Var(\phi_j)$
using i.i.d.\ seeds and the covariance inequality
(features within the same group in the same model are positively
correlated, but the bound uses only the marginal variance, giving a
conservative upper bound).
Concentration in probability gives convergence to distinct
limits across groups, ensuring the group-level rank order stabilises.

Part~(c) follows from the symmetry of~(a): $\bar\phi_j - \bar\phi_k \sim
N(0,\,\Theta(1/M))$ by CLT (independent seeds, bounded variance), and two
independent runs produce independent draws of this quantity; the probability
they share the same sign is $1/2 + O(1/\sqrt{M})$, not tending to~$1$.
\end{proof}

\begin{remark}[Within-group instability is not a pathology]
\label{rem:within-group}
\Cref{cor:dash}(c) may appear to concede ground---\dash{} achieves equity
but leaves within-group rankings noise-dominated.  However, this is the
\emph{correct} behaviour when features are genuinely equivalent in the DGP:
any attribution method that consistently ranks $j$ above $k$ when
$\beta_j = \beta_k$ is introducing a spurious asymmetry.  The relevant
interpretability target is identifying \emph{which group} is important
(between-group stability) and confirming that no individual feature within
the group is systematically advantaged (equity).  Both are achieved.
Individual within-group rankings are interpretively meaningless when
$\beta_j = \beta_k$, and the appropriate diagnostic for detecting
within-group instability is the FSI (Feature Stability Index, defined in
the main paper), not rank correlation.
\end{remark}


% ══════════════════════════════════════════════════════════════════════════════
\section{Discussion of Assumptions}
\label{sec:assumptions}
% ══════════════════════════════════════════════════════════════════════════════

\begin{remark}[Sequential structure is essential]
\label{rem:sequential}
The impossibility in \Cref{thm:impossibility} relies critically on the
\emph{sequential} nature of gradient boosting, where each tree $h_t$ is
fit to residuals $r_t = y - \sum_{s<t} h_s(x)$.  This creates a causal
chain: the first tree's split choices alter the residual landscape for all
subsequent trees, producing the cumulative first-mover advantage of
\Cref{lem:split-concentration}.

The theorem does \emph{not} apply to methods that train base learners
independently.  For example:
\begin{itemize}
  \item \textbf{Random forests} train each tree on a bootstrap sample with
    independent feature sub-sampling.  Since trees do not share residuals,
    there is no first-mover propagation: if tree $t$ splits on $j$, this
    does not affect subsequent trees' feature selection.  The equity violation
    is therefore $O(1)$ rather than $O(\rhovar^2 \cdot \totalT)$.
  \item \textbf{Bagged ensembles} similarly avoid the sequential dependency
    by construction.
\end{itemize}
However, gradient boosting's sequential structure is precisely what makes
it a powerful function approximator (by fitting residuals, it can represent
complex interactions that parallel ensembles miss).  The impossibility
result thus highlights a fundamental tension between the \emph{modeling
power} of sequential boosting and the \emph{fairness properties} of
feature attribution.
\end{remark}

\begin{remark}[Scope of \Cref{ass:uniform}]
\label{rem:assumption-scope}
\Cref{ass:uniform} (the uniform-contribution model) abstracts away the
depth-weighting of TreeSHAP, which assigns higher weight to splits closer
to the root.  In practice, early splits on the first-mover feature
$j^*(\omega)$ occur preferentially near the root (since the first tree's
root split is on $j^*(\omega)$), so the actual attribution gap is likely
\emph{larger} than the $\Omega(\rhovar^2)$ established here.  The
assumption therefore gives a conservative lower bound.  A tighter analysis
tracking depth-weighted split counts would strengthen the impossibility
result.
\end{remark}

\end{document}
